\subsection{Mixed first coordinates at three and a split prime}

\begin{statement}{E.56}{Proved}{Every mixed first-coordinate fiber at an explicitly specified split prime}
Let $\ell$ be a rational prime satisfying
\[
       v_3(\ell-1)=1,\qquad
                 3^{(\ell-1)/3}\not\equiv1\pmod\ell .
 \tag{E56.1}
\]
For example, $\ell=7,13,31,43,79,97$ satisfy these conditions.
Use the actual arithmetic quotient $H$ and the scaled toric
map of E.46. Fix $a\ne0$ and $b$ in $\mathbb Z_3$, and put
\[
 \begin{gathered}
 A=3a=3^s\epsilon,\quad s\ge1,\quad\epsilon\in\mathbb Z_3^\times,
 \qquad B=3b,\qquad h=(\ell-1)/3\in\mathbb Z_3^\times,\\
 r=B/(\ell-1)=b/h\in\mathbb Z_3,\qquad
 [\mathcal A]=[3a\kappa_3(\ell)+3b\kappa_3(3)].
 \end{gathered}
 \tag{E56.2}
\]
For $p\notin\{\ell,3\}$ let $h_p=v_3(p-1)$, choose
integer representatives $a_p,b_p$ modulo
$3^{\max(h_p-1,0)}$, and define
\[
 \mathscr P=
 \left\{p\notin\{\ell,3\}:
       \bigl(\ell^{3a_p}3^{3b_p}\bigr)^{(p-1)/3^{h_p}}
                                     \equiv1\pmod p\right\}.
 \tag{E56.3}
\]
The condition is independent of the representatives.
Let $\mathcal R(\mathscr P)$ have the restricted
inverse-limit topology of (E49.2).

All global section classes of $H$ with first class
$[\mathcal A]$ have complete additive topological coordinates
\[
 \boxed{\qquad
 \mathbb Z_3^2\times\mathcal R(\mathscr P)
                    \times\mathbb Z_3^{\,2s}\times\mathbb Z_3[[Z]].
 \qquad}
 \tag{E56.4}
\]
The first two coordinates are denoted $\beta_\ell,\beta_3$,
and the restricted ones $\beta_p$. The scaled toric map is
exactly
\[
 \left(a\kappa_3(\ell)+b\kappa_3(3),\
   \beta_\ell\kappa_3(\ell)+\beta_3\kappa_3(3)
                      +\sum_{p\in\mathscr P}\beta_p\kappa_3(p)\right).
 \tag{E56.5}
\]
Every allowed second coordinate has the explicitly
computed compact fiber
$\mathbb Z_3^{2s}\times\mathbb Z_3[[Z]]$.
The entire group in (E56.4), with the zero second
cocycle as origin, is torsion-free.
The first classes with $a=0$ are already computed
in E.47 and E.51.

Here are actual arithmetic generators and the complete
normal form. Put $K=\mathbb Q(\zeta_3)$ and let $S$
contain the primes over $\ell$ and three. There are
actual local generators $t,\sigma$ at one prime over
$\ell$, and $t'=ct^{-1}c$, $\sigma'=c\sigma c$ at its
conjugate, such that
\[
 \begin{gathered}
 G_{K,S}^{(3)}
   \simeq
   \langle t,\sigma\mid\sigma t\sigma^{-1}=t^\ell\rangle_3
       \mathop{\ast}_3
   \langle t',\sigma'\mid\sigma't'\sigma'^{-1}=t'^\ell\rangle_3,\\[2mm]
 \begin{array}{c|rrrr}
       &t&t'&\sigma&\sigma'\\\hline
 \chi_3&1&1&\ell&\ell\\
 \kappa_3(\ell)&1&1&0&0\\
 \kappa_3(3)&0&0&1&-1 .
 \end{array}
 \end{gathered}
 \tag{E56.6}
\]
The full relative arithmetic group over $C_2$
projects onto this group by a map with an actual
$c$-equivariant section retaining both decomposition
groups.

For $1\le d\le s$ and $\delta\in\{+1,-1\}$ define
the centered orbit sums
\[
 Q_{d,\delta}(T)=
   \sum_{\substack{-(3^d-1)/2\le j\le(3^d-1)/2\\
                                  j\equiv\delta\pmod3}}
                      T^{\epsilon3^{s-d}j}.
 \tag{E56.7}
\]
The finite parameter lattice has the explicit basis
given by the free coordinates
\[
  \eta,\ z_{1,-},\
                  (z_{d,+},z_{d,-})_{2\le d\le s}\in\mathbb Z_3,
 \qquad
 z_{1,+}=3\eta+\sum_{d=1}^s z_{d,-}
                           -\sum_{d=2}^s z_{d,+}.
 \tag{E56.8}
\]
There are exactly $2s$ free coordinates.
Put
\[
 \begin{split}
 f_0(T)&=3\beta_\ell+
       \sum_{d=1}^s\sum_{\delta=\pm1}
          z_{d,\delta}\bigl(Q_{d,\delta}(T)-3^{d-1}\bigr),\\
 u_0(T)&=T^{-r}f_0(T),\qquad
 S_\ell(T)=\sum_{i=0}^{\ell-1}T^{iA},\\
 v_0(T)&=
       \frac{\ell T^B u_0(T^\ell)-S_\ell(T)u_0(T)}
                                                {T^{\ell A}-1},\\
 V&=T+T^{-1}-2,\qquad
 z(T)=\frac{\beta_3}{h}-\eta+\frac{3\beta_\ell}{2}
                                                   +Vf(V).
 \end{split}
 \tag{E56.9}
\]
The quotient defining $v_0$ is integral. For an arbitrary
series $f\in\mathbb Z_3[[Z]]$, the unique normal representative
of the class unramified outside $S$ is
\[
 \begin{split}
 v(c)&=0,\\
 u=v(t)&=u_0+(T^A-1)z,\\
 v(\sigma)&=v_0+(\ell T^B\psi_\ell-1)z,
                       \qquad\psi_\ell(g)=g(T^\ell),\\
 v(t')&=T^A u(T^{-1}),\qquad
 v(\sigma')=-v(\sigma)(T^{-1}).
 \end{split}
 \tag{E56.10}
\]
Its augmentations on $t,\sigma$ are exactly
$3\beta_\ell,3\beta_3$.
The proof constructs the additional global residue
cocycles and their exact toric normalization.
It also gives the forward and inverse maps to the
actual norm systems in E.46, rather than identifying
the parameter products without a comparison map.

The maps from the intrinsic Fermat groups to this
$H$ retain the kernels displayed in (E47.8).
Equations (E56.4)--(E56.10) evaluate the target
fibers for the stated first-coordinate family.
They leave the source lifting conditions attached
to those exact kernels. The manuscript passage is
\S51.16, folios~733v--736v, scans~172--178,
equations~(371)--(379), Part E(c),(d).
\end{statement}

\begin{proof}
\emph{The two actual decomposition factors.}
Choose compatible roots of the positive rational numbers
fixed by complex conjugation $c$. Since $\ell\equiv1\pmod3$,
it splits in $K$. Choose the prime at which the cubic
Kummer Frobenius value of three is one modulo three.
Such a choice exists by (E56.1); conjugation changes
the value to its negative. Choose an actual Frobenius
$\sigma_0$ fixing the compatible roots of the
uniformizer $\ell$, and put $\sigma'_0=c\sigma_0c$.
Write $t_0=\kappa_3(3)(\sigma_0)\equiv1\pmod3$.
The exact Kummer values are
\[
 \chi(\sigma_0)=\chi(\sigma'_0)=\ell,\qquad
 \kappa_\ell(\sigma_0)=\kappa_\ell(\sigma'_0)=0,\qquad
 \kappa_3(\sigma_0)=t_0,\quad
 \kappa_3(\sigma'_0)=-t_0 .
 \tag{E56.11}
\]
Here $\kappa_\ell=\kappa_3(\ell)$ and
$\kappa_3=\kappa_3(3)$ in this calculation.
The cyclotomic character modulo nine and
$\kappa_3$ modulo three give the two independent
Frattini characters of the free rank-two group
$P=G_{K,\{3\}}^{(3)}$ proved in E.34.
The two Frobenii have determinant
$-2((\ell-1)/3)t_0$, a three-adic unit, in these
coordinates. Their images are therefore a free
basis of $P$.

An explicit conjugation makes the displayed values exact,
not just values modulo three. Put
\[
 \rho_0=\sigma'_0\sigma_0^{-1},\qquad
 k=\frac{1-t_0}{2t_0(\ell-1)}\in\mathbb Z_3,\qquad
 g=\rho_0^k,\qquad
 \sigma=g\sigma_0g^{-1},\quad \sigma'=c\sigma c .
 \tag{E56.12}
\]
The integrality of $k$ follows from
$t_0\equiv1\pmod3$ and $v_3(\ell-1)=1$.
Since $\chi(\rho_0)=1$,
$\kappa_\ell(\rho_0)=0$, and
$\kappa_3(\rho_0)=-2t_0$, the cocycle identity gives
\[
 \kappa_3(\sigma)
   =t_0+2t_0(\ell-1)k=1.
 \tag{E56.13}
\]
All other entries in (E56.11) are preserved.
Conjugating each of a free pair does not change
its Frattini images, so the new images remain
a free basis. Conjugate the first tame inertia
generator by $g$ as well, normalize its
uniformizer Kummer value to one, and call it $t$.
Set $t'=ct^{-1}c$. This proves the entire table
(E56.6), with actual local generators.

At each prime the local pro-three tame group
injects into the full pro-three arithmetic group:
its inertia injects by the free-inertia theorem,
and its unramified part injects through the
nontrivial cyclotomic character. Applying the
opened Wingberg Theorem~2.4
\cite{E_WingbergRiemann} to $\{3\}\subset S$
identifies the kernel over $P$ with the continuous
free inertia bundle over the two coset spaces
$P/\langle\sigma\rangle$ and
$P/\langle\sigma'\rangle$.
The $K(\pi,1)$ hypothesis and nontrivial local
unramified completions are the ones verified
in E.39 and E.55.

There is an explicit inverse for the resulting
free-product comparison. For either generator
$\tau=\sigma,\sigma'$, form the induced
inertia bundle with
$(p\tau^z,t_\tau^u)$ identified with
$(p,t_\tau^{\ell^zu})$.
Take the continuous free pro-three product of
these two bundles and its semidirect product
with $P=\langle\sigma\rangle\ast_3\langle\sigma'\rangle$.
The inclusions of $t,\sigma,t',\sigma'$ give maps
in both directions between this semidirect
product and the right side of (E56.6).
They are inverse on those generators, by the
free-product universal property.
On the arithmetic kernel, Wingberg's
isomorphism sends every factor to its actual
inertia group, and the map on $P$ is the
identity. This proves the isomorphism
(E56.6) including its kernel map.
We chose the generators in the full
arithmetic group, so their same inclusions
define a reverse section of its projection
onto (E56.6). Its projection fixes every
generator, and it commutes with $c$.
Enlarging $S$ to all primes in the same
theorem identifies the remaining kernel
as the continuous free inertia product
outside $S$.

\emph{The exact two-generator cocycle calculation.}
For the first cocycle $\mathcal A$, the coefficient
module is $J=(3,T-1)\Lambda$, with
$\Lambda=\mathbb Z_3[[T-1]]$.
On the first decomposition group its operators are
\[
    t(f)=T^Af,\qquad
    \sigma(f)=\ell T^B f(T^\ell),\qquad
    C(f)=-f(T^{-1}).
 \tag{E56.14}
\]
Normalize the second cocycle so that $v(c)=0$,
using division by two. Since $c$ exchanges
the decomposition factors, its values on the
second factor are then exactly the last line
of (E56.10). There are no further relations
between the two factors. The single condition
on $u=v(t)$ and $v=v(\sigma)$ is
\[
       \ell T^Bu(T^\ell)-S_\ell(T)u(T)
                                      =(T^{\ell A}-1)v(T).
 \tag{E56.15}
\]
This follows directly from the cocycle identity
for $\sigma t\sigma^{-1}=t^\ell$, and is
sufficient by the actual presentation.

The remaining geometric conjugations have
zero first coordinate: a first coordinate
preserving the value at $c$ is killed by
multiplication by two. Their second
coordinates are exactly the series
$m$ with $m(T^{-1})=-m(T)$.
They act by
\[
       (u,v)\longmapsto
       \bigl(u-(T^A-1)m,\ 
             v-(\ell T^B\psi_\ell-1)m\bigr).
 \tag{E56.16}
\]
Every such $m$ has augmentation zero
and hence belongs to $J$.

Put $M=3^s$. The ideals
$(T^A-1)$ and $(T^{\ell A}-1)$
both equal $(T^M-1)$.
Their ratios are power series with
constant terms $\epsilon$ and $\ell$,
respectively, and are units in $\Lambda$.
Thus the quotient is the actual cyclic
group ring $\mathbb Z_3[C_M]$.
Because $B=(\ell-1)r$, multiplication
by $T^r$ turns the affine operator
$T^B\psi_\ell$ into $\psi_\ell$.
Reducing (E56.15) consequently says
exactly that
\[
             \bar f=\overline{T^r u}
              \quad\text{is fixed by multiplication by }\ell
                        \text{ in }\mathbb Z_3[C_M].
 \tag{E56.17}
\]
The closed subgroup generated by $\ell$
is $1+3\mathbb Z_3$. At each nonzero
order $3^d$ there are precisely two
orbits, represented by
$\epsilon3^{s-d}$ and
$-\epsilon3^{s-d}$. Their orbit sums
are the images of (E56.7).
Together with zero, these are the
$2s+1$ orbits of the complete
finite permutation module.

Augmentation of $u$ and $f$ agrees.
Thus all invariant elements with
augmentation in $3\mathbb Z_3$ have
the unique lift $f_0$ in (E56.9),
with $\beta_\ell$ and all
$z_{d,\delta}$ initially arbitrary.
The numerator defining $v_0$
vanishes modulo $T^A-1$ by
(E56.17), and therefore is divisible
by $T^{\ell A}-1$ in $\Lambda$.
This proves the claimed integral
division.

There is one additional, exactly
computable condition for $v_0$ to
belong to $J$. Let
$\partial=T\,d/dT$ and evaluate
at $T=1$. The centered representatives
satisfy
\[
 \begin{aligned}
 \operatorname{aug}Q_{d,\delta}&=3^{d-1},&
 (\partial Q_{d,\delta})(1)&=\delta A/3,\\
 \operatorname{aug}v_0
   &=\frac{\ell-1}{A}
       \left((\partial f_0)(1)
                       -\frac A2\operatorname{aug}f_0\right)\\
   &=(\ell-1)\left(
          \frac13\sum_{d=1}^s(z_{d,+}-z_{d,-})
                                      -\frac{3\beta_\ell}{2}\right).
 \end{aligned}
 \tag{E56.18}
\]
For the first-moment identity, the
$3^{d-1}$ centered integers congruent
to $+1$ modulo three have mean one;
the negative orbit has mean minus one.
The derivative formula follows by
differentiating the actual numerator
and denominator of (E56.9); the
$r$ terms cancel because
$B=(\ell-1)r$.
As $(\ell-1)/3$ is a unit,
$\operatorname{aug}v_0\in3\mathbb Z_3$
is equivalent to
\[
          \sum_{d=1}^s(z_{d,+}-z_{d,-})\in3\mathbb Z_3.
 \tag{E56.19}
\]
This is exactly the lattice with
the explicit basis (E56.8).
It has index three in the
free $2s$-coordinate lattice;
the stated formula solves that
condition integrally rather than
dividing a free coordinate by
three.

Every solution with these finite
coordinates has a unique expression
\[
       u=u_0+(T^A-1)z,\qquad
       v=v_0+(\ell T^B\psi_\ell-1)z,
                              \qquad z\in\Lambda.
 \tag{E56.20}
\]
Indeed the first difference is
divisible by $T^A-1$, and the
second follows from (E56.15)
in the integral domain $\Lambda$.
If an arbitrary solution starts
with $u,v\in J$, then
$\operatorname{aug}v-
\operatorname{aug}v_0
=(\ell-1)\operatorname{aug}z$
is already divisible by three.
It follows that (E56.19) is
necessary for every solution,
not only for the proposed lift.

Split $z$ into its even and
odd parts under $T\mapsto T^{-1}$.
The unique allowed conjugation
in (E56.16) removes its odd part.
The even ring is
$\mathbb Z_3[[V]]$:
with $W=(T-T^{-1})/2$,
the ring is $\mathbb Z_3[[W]]$,
inversion sends $W$ to $-W$,
and
$V=2(\sqrt{1+W^2}-1)$ has
unit coefficient of $W^2$.
The desired augmentation
$\operatorname{aug}v=3\beta_3$,
together with (E56.18) and
$\sum(z_+-z_-)=3\eta$,
fixes the constant part of $z$
to be exactly
\[
             \operatorname{aug}z
                =\beta_3/h-\eta+3\beta_\ell/2.
 \tag{E56.21}
\]
Its remaining even part is
uniquely $Vf(V)$.
This proves every formula in
(E56.9)--(E56.10), its
existence and its uniqueness.
No further nonzero geometric
conjugation fixes $u$, since
$(T^A-1)m=0$ in a domain
forces $m=0$.
The specified substitutions,
finite coefficient extraction,
division on closed principal
ideals, and even/odd projectors
are continuous. This proves the
full compact normal-form
homeomorphism.

\emph{The entire global section fiber.}
Let $L$ be the kernel of the
actual split arithmetic
projection onto (E56.6).
Its action on $J_{\mathcal A}$
is trivial. As in the exact
reconstruction (E55.13), a
global second cocycle is
equivalent to its cocycle on
(E56.6) and an equivariant
continuous additive homomorphism
$\lambda:L\to J_{\mathcal A}$.
The free inertia bundle
computes these homomorphisms.
At $p\notin\{\ell,3\}$ their
possible normalized inertia
values are zero or
$3\mathbb Z_3\delta_{x_p}$,
where
$x_p=-\mathcal A(F_p)/(p-1)$.
The fixed-measure proof in
E.49 and E.55 shows that this
Dirac line occurs exactly when
$x_p\in\mathbb Z_3$.
For the actual Kummer cocycle
this is precisely the
three-primary residue
condition (E56.3).
At a split prime its
conjugate value is determined
by $c$; at an inert prime
the rational Frobenius
condition supplies the
same descent. The family
of coefficients is exactly
the restricted sequence
module $\mathcal R(\mathscr P)$.

For clarity, the rational
toric normalization differs
from the one used for the
inert decomposition factor
in E.55 by an explicitly
computed evaluation map.
Choose roots of every
rational prime fixed by
$c$. A residue homomorphism
supported at $p$, with
inertia value $3\delta_{x_p}$,
and zero cocycle on the
split arithmetic subgroup,
has scaled second Kummer
class exactly
\[
                    \kappa_3(p)
                       -\kappa_3(p)(\sigma)\kappa_3(3).
 \tag{E56.22}
\]
Its valuation is one at
$p$, zero at every other
prime except possibly
three, and zero at $\ell$.
It vanishes on $\sigma$
and $\sigma'$; the
two evaluation identities
$\kappa_3(3)(\sigma)=1$
and
$\kappa_3(p)(\sigma')
=-\kappa_3(p)(\sigma)$
fix the coefficient at
three exactly as displayed.
These facts determine
the rational completed
Kummer class by E.16.

The normal form
(E56.10) with
$\beta_\ell=0$,
all finite coordinates
and $f$ zero, and
$\beta_3=\kappa_3(p)(\sigma)$
has scaled toric class
$\kappa_3(p)(\sigma)\kappa_3(3)$.
Adding it to the residue
cocycle produces an
actual global cocycle
$v_p$ with toric class
$\kappa_3(p)$.
Every global cocycle has
unique residues
$3\beta_p\delta_{x_p}$.
Subtracting
$\sum_p\beta_pv_p$
leaves precisely a
cocycle on the subgroup
(E56.6), for which the
normal form is unique.
The series converges
in every finite
coefficient and
positional quotient:
$\beta_p\to0$, and the
normalizing coefficients
$\kappa_3(p)(\sigma)$
are integral.
Valuation extraction,
the normal form, and
this reconstruction
are continuous in the
uniform cochain
quotient topology.
This proves
(E56.4)--(E56.5) as
a full global
homeomorphism.

\emph{Comparison with the norm calculation and verification.}
The map to E.46 is fully
specified. Starting with
a norm-compatible system,
apply its finite Kummer--Shapiro
inverse, align the compatible
cocycles as proved there,
and fix the first representative
$\mathcal A$. Normalize
$v(c)=0$. Extract the
fixed inertia measures
at primes outside $S$
and subtract the just
specified $v_p$ terms.
For the remaining
cocycle, take
$\overline{T^r v(t)}$
in $\mathbb Z_3[C_{3^s}]$
and read its orbit
coefficients. Divide
$v(t)-u_0$ by $T^A-1$,
project to its even
part, and use
(E56.21) to extract
the final power series.
This is the forward
map to (E56.4).
Conversely the actual
arithmetic presentation
constructs the cocycle
(E56.10), the additive
inertia reconstruction
adds the global $v_p$
terms, and E.46 sends
it to the norm system
by the stated finite
Kummer--Shapiro maps.
The equations above
and the inverse theorem
in E.46 prove that
these two maps are
mutual inverses at
every finite level
and hence in the
stated topologies.

The retained exact
program
\texttt{e\_split\_prime\_first.py}
checks the actual
Kummer normalization
and the Frattini
determinants, the
complete finite orbit
systems, both
first-moment identities,
the index-three lattice,
the integral Laurent
division and its
unit-denominator
series, and recovery
of every tested
finite, toric, odd,
and even coordinate.
It also checks the
eligible-prime
criterion by four
finite descriptions.
The proof establishes
the infinite arithmetic
and every three-adic
parameter; the finite
certificates verify
the explicit formulas
and their signs.
\end{proof}
